\gddchapter{Interview sequence}


Below is the guideline for the semi-open guided interview. 
The interview starts with easy warm-up questions and progresses towards harder questions as it goes along. 

The answers were translated from polish back to english.

\section{Introduction}
Remind the participant that this is not a test, rather a joint exploration.

\section{Icebreakers}

\begin{QandA}
   \item Now that you’ve tried both versions, how are you feeling? Was there anything that immediately surprised you?
         \begin{answered}
            The first thing is that the environment is described really well, and that makes the images not that necessary, because it shows that you can imagine the space spatially.

            Only that by seeing the options faster, I could skip some of the descriptions, since I already see what the options are and so on, but that’s not really the point in the end.

            And the fact that you actually immerse yourself-when you have to give a command, you need to know what’s happening around you.

            So for example, if there were no text where you highlight those objects, mark those paths, it could be a bit confusing, because you wouldn’t really know what you can do or what you can pick up.

            It would have to be a really… I’d have to just arrive at a fully interactive world and environment. 

            But..

            I was actually wondering how to phrase commands so that it would definitely understand me.
            When it comes to being surprised there was no surpirses. It was just exploring.
                  \end{answered}

   \item If you had to describe the difference between these two games to a person who hasn't played either, what would you say?
         \begin{answered}
            It’s like I’m reading a book and have an influence over what happens next, choosing options.

            I’m kind of reading a book, but also talking to it.

            I also wasn’t sure-I was trying to test it, like if I wasn’t sure whether I could go somewhere, I asked a question: “can we go?”

            And maybe it would be nice if I got a more complete answer, like “no, because the door is locked,” for example.

            Or that you can’t pick something up that’s mentioned in the text because it’s too heavy.

            That would create this feeling that you have to think and actually read the text carefully.

            But then the text would have to be refined to that level.

            Or maybe it should be kind of vague, so you don’t really know—that could also be part of the fun.

            That you test things and get different responses, right?

            That’s an idea, maybe.

            I just didn’t really know what I could do.

            With the options, it felt clearer.
         \end{answered}
\end{QandA}




\section{Grand tour questions}
(remember to pause here 3-5 seconds)

\begin{QandA}


   \item From the moment you first used your voice to navigate, what were you thinking as you decided what to say to the system?

         \begin{answered}
            Simple commands.

            The same ones as before.

            I tried to use them the same way they were written earlier.

            So like “go to,” “pick up.”

            I would probably also use voice.

            When I wasn’t sure - whether I was blocked or could go through - I asked things like “can we go,” because we’re also aware of the limitations.
         \end{answered}

    \item How was it like to use the voice input in the game?
         \begin{answered}
            Because those commands were short - I was influenced by the option selection - it felt like I wanted to type something into a search engine.

            So it’s like I don’t have to click, but I’m kind of used to playing like that.

            So it’s about breaking away from what I’m used to when playing a game.

            I was more immersed in the second one.

            The first one allowed for quick skipping.

            If I played longer, speaking would probably become easier, because there were a lot of transitions like “travel to,” “go back to main list,” and “use item” that you had to choose.

            Actually, if I played longer using voice, I would only have the inventory and that list while reading.

            It would be faster to speak.

            But I’d need to play longer to get more familiar with how it works.
         \end{answered}


\end{QandA}




\section{Core comparative questions}

\begin{QandA}
   \item In which version did you feel more 'inside' the story world?
         \begin{answered}
            The second one - that’s more like immersing yourself in the world. Deeper.”
            “I would choose the second one.
         \end{answered}

   \item How did the 'effort' of thinking of what to say naturally compare to the 'effort' of knowing which keys to press?

         \begin{answered}
            But because I didn’t really know how much freedom I had when speaking, then…

            If it were interactive in a way where I could ask questions and get answers tied to the game - like I said, asking “can we go there?” and getting “no, because…” - and it would pull, for example, a fragment from the text and the reason why we can’t go there…

            That would be great.
         \end{answered}

    \item Did the freedom to say 'anything' make you feel more in control, or was it more overwhelming than having a fixed list of keys?

         \begin{answered}
            The first one - boring.

            Simple, fast, just to see the version of the world you imagined in the apocalypse, to get familiar with it.

            Quick.

            And the second one?

            That’s more like immersing yourself in the world.

            Deeper.

            Deeper.

            And if the dialogues were even more developed?

            That would be… powerful.
         
         \end{answered}
    \item How would you compare the overall experience of using the two game versions? Which one would you choose and why?

         \begin{answered}
         
         \end{answered}
\end{QandA}






\section{Probing}
In this part of the interview look into the sticky moments noted above and ask follow-up questions probing for more detail.
Include lexical reflection (eg. you used the word "clunky" when explaining navigation. Could you tell me more about what was happening
in that moment?)


\section{Final reflection}

\begin{QandA}
   \item Okay, and from what you're saying, it sounds like this graphic could also play an additional role in the experience to some extent.
         \begin{answered}
         Yes, but wouldn’t that dilute the experience? Because all… right, that’s another issue. Are you the type who prefers movies, or the type who prefers books? You know, there are people… which type are you? I like books—I like books for a different reason. I’m aware of the difference in experiences, so I choose which kind of experience I want to have. A movie is fast, while a book is engaging; when I read long books, I sometimes have to take breaks so I don’t get too absorbed in the story, because I experience it too intensely and connect with the characters so deeply. They exist only in my head, so I imagine them however I want and feel them. You also can’t hear thoughts just by watching an actor cry—if something isn’t said, you won’t know it—but in a book, you can know the characters’ thoughts and so on. So these are two completely different experiences. Film—images—are much more shallow than what you can achieve with text. It’s hard to compare them, because they’re based on completely different intentions—just a different kind of experience.
         \end{answered}

   \item After reflecting on everything today, is there anything else you’d like to add about using your voice to play this game that we haven't talked about?

         \begin{answered}
         Not for now, but maybe I’ll dream about it and then I’ll message you.
         \end{answered}
\end{QandA}

\section{Closing}
Thank the participant for his time and tell them again how much value his input and time gives me.






